\documentclass{situationist-preprint}

\title{Omniscience Without Tensed Belief}
\msid{REL-19-b}
\msfamily{REL-19}
\msdiscipline{Philosophy of Religion}
\mstype{Consequence}
\mskeywords{omniscience; timelessness; tensed belief; indexicals; self-locating belief; Kretzmann}

\begin{document}
\maketitle

\begin{abstract}
REL-19-a concludes that God is timeless and therefore lacks indexical knowledge of what time it is now, and that this is no defect in omniscience properly defined. We derive three consequences. A timeless God has no tensed beliefs at all. This leaves no fact about the world unknown to God, since on a tenseless view of time \enquote{now} carries self-locating information and not information about which world is actual. And God needs no tensed belief in order to act at the right time, because God's will concerns events at times and is not exercised from a time. Kretzmann's argument that an omniscient being must change, and Prior's argument from relief, therefore leave timeless omniscience untouched.
\end{abstract}

\section{No tensed beliefs}

A \emph{tensed belief} is a belief, held at a time, in a centered proposition about the believer's temporal position: that it is now noon, that the meeting is starting, that the pain is over. In the terms of REL-19-a (Definition 2.1), its content is a set of centers \citep{Lewis1979}, and it is held de nunc.

\begin{proposition}\label{p:notensed}
God has no tensed beliefs, and God's knowledge consists entirely of tenseless propositions.
\end{proposition}

\begin{proof}
A tensed belief is held at a time, and believing de nunc, like knowing de nunc, presupposes a location at that time (Assumption 2.5 of REL-19-a and the account of Definition 2.3 behind it). By Theorem 3.3 of REL-19-a, God is located at no time. By Corollary 4.2, God knows no centered proposition, and by Assumption 4.1, God knows every true tenseless proposition.
\end{proof}

\section{Nothing about the world is missed}

TSM-01 concludes that time is real and constituted by the B-series: the temporal facts are facts of earlier, later and simultaneous, and there is no further fact of which time is present. TSM-18 adds that \enquote{now} is an indexical, so that there is no standpoint-free fact of what time it is now. EPI-18-c applies the same view to self-locating belief: a subject's uncertainty about whether it is now Monday is uncertainty about where she is located in a world, and learning that it is Monday is conditionalization on self-locating information, not the acquisition of a tensed fact.

\begin{proposition}\label{p:world}
Any two possible worlds that differ in some respect are distinguished by a true tenseless proposition that God knows.
\end{proposition}

\begin{proof}
On Definition 2.1 of REL-19-a, a tenseless proposition is a set of worlds, and for any two distinct worlds there is a set containing one and not the other; the one containing the actual world is true and, by Assumption 4.1, known to God. Centered propositions add no discrimination among worlds beyond this: they discriminate among centers within a world, and on the B-theory of TSM-01 a world contains no further fact of which center is present.
\end{proof}

Proposition~\ref{p:world} is the sharp form of Proposition 4.4 of REL-19-a, which shows that whatever a temporal subject knows de nunc, God knows the corresponding tenseless truth. What God lacks is not a fact but a location. Ignorance in the defective sense is failing to know how things are; a timeless God knows how things are throughout the B-series, including that at $t$ a subject knows that it is then $t$.

\section{Acting at the right time}

\citet{Perry1979} showed that indexical belief is essential to timely action. The professor who knows that the meeting starts at noon does not leave until he believes that it is noon now. It may seem that a God without tensed belief could not act at the right time.

\begin{proposition}\label{p:action}
God's acting so that an event occurs at $t$ requires no tensed belief.
\end{proposition}

\begin{proof}
A temporal agent needs a tensed belief because she initiates action from her own location and must locate the moment of action relative to it. God's will, being timeless, is not exercised from a location: its content is tenseless, that event $e$ occur at $t$, and its effect is located at $t$ without being initiated from an earlier or later time \citep{StumpKretzmann1981}. The role that tensed belief plays for a temporal agent, linking a tenseless plan to a moment of initiation, has no application.
\end{proof}

\section{Two objections}

\begin{objection}
\citet{Kretzmann1966} argued that a perfect knower knows what time it is, that what time it is keeps changing, and that such a knower therefore changes. If omniscience requires knowing what time it is, a timeless God is not omniscient; if it does not, something is unknown to God.
\end{objection}

\begin{reply}
The second horn assumes that what time it is is something to be known beyond the B-series facts. On the view of TSM-01 and TSM-18 it is not: \enquote{what time it is} is fixed only relative to a center. By Definition 2.6 of REL-19-a, omniscience requires knowledge of centered truths only at the knower's own center, and Proposition 4.5 shows that the stronger requirement, knowledge of every centered truth, can be met by no being at all. \citet{Grim1985} takes that impossibility to tell against omniscience; it tells instead against the stronger definition.
\end{reply}

\begin{objection}
\citet{Prior1959} observed that, after a painful experience, one is relieved that it is over, and that what one is relieved about is not a tenseless fact. A God who knows only tenseless facts cannot share in creatures' relief, or rejoice that their suffering has ended.
\end{objection}

\begin{reply}
Relief is an attitude of a subject located after the pain toward its own temporal position, and its content is centered, as EPI-18-c's treatment of \enquote{now} predicts. A being located at no time has no position after the pain from which to be relieved. What God knows, tenselessly, is that the suffering ends at $t$ and that the creature is relieved after $t$, and God's response to that is not deficient for being unindexed.
\end{reply}

\section{Conclusion}

A timeless God believes nothing in the tensed mode, and loses nothing by it. The world contains no tensed facts for God to miss, God's acts need no indexical trigger, and the arguments of Kretzmann and Prior show only that a timeless God has no location in time, which is what REL-19-a concludes.

\begin{thebibliography}{99}\small
\bibitem[Grim(1985)]{Grim1985} Grim, P. (1985). Against omniscience: The case from essential indexicals. \emph{No\^us} 19(2): 151--180.
\bibitem[Kretzmann(1966)]{Kretzmann1966} Kretzmann, N. (1966). Omniscience and immutability. \emph{Journal of Philosophy} 63(14): 409--421.
\bibitem[Lewis(1979)]{Lewis1979} Lewis, D. (1979). Attitudes de dicto and de se. \emph{Philosophical Review} 88(4): 513--543.
\bibitem[Perry(1979)]{Perry1979} Perry, J. (1979). The problem of the essential indexical. \emph{No\^us} 13(1): 3--21.
\bibitem[Prior(1959)]{Prior1959} Prior, A.~N. (1959). Thank goodness that's over. \emph{Philosophy} 34(128): 12--17.
\bibitem[Stump and Kretzmann(1981)]{StumpKretzmann1981} Stump, E. and Kretzmann, N. (1981). Eternity. \emph{Journal of Philosophy} 78(8): 429--458.
\end{thebibliography}

\end{document}
